\chapter{Background}
\label{chap:background}

This chapter establishes the conceptual foundations that underpin the rest of the thesis. The chapter first frames stock forecasting as a problem of limited predictability, then introduces target definitions and input modalities, the role of class imbalance, the use of textual information in finance, the move from lexicons to deep financial Natural Language Processing (NLP), Arabic and bilingual financial text processing, and chronological evaluation with explicit leakage control. Together, these concepts define the methodological surface on which Chapter~\ref{chap:methodology} operates.

\section{Financial Forecasting as a Problem of Limited Predictability}

Stock forecasting sits at the intersection of economics, statistics, and machine learning. Its difficulty begins with the nature of the market itself. Financial prices respond to heterogeneous information, strategic trading, liquidity effects, and macroeconomic shocks. The efficient-market literature formalised this challenge by arguing that public information is rapidly incorporated into prices. Later empirical work showed that the random-walk view is too strict to hold universally, but it also confirmed that predictable structure in market data is weak, unstable, and difficult to exploit consistently \cite{fama1970, lo1988randomwalk}.

This tension explains why forecasting remains an active research area. If markets were perfectly predictable, the problem would already be solved. If they were perfectly unpredictable, forecasting models would be pointless. The research problem lies between these extremes. Small predictive patterns may exist, but they are noisy, regime-dependent, and sensitive to how the prediction target is defined. Surveys of the field show that forecasting studies differ substantially in horizon, label construction, feature design, and evaluation protocol. Those choices often matter as much as the choice of learning algorithm itself \cite{atsalakis2009survey, sezer2020dlreview}.

\section{Forecast Targets and Input Modalities}

Financial forecasting studies use several related targets. Some studies predict the next price level. Others predict return magnitude, return sign, or discrete upward versus downward direction. Directional forecasting aligns naturally with decision-making, but it discards part of the information contained in the return magnitude. Regression targets preserve more information yet are often harder to estimate accurately. Many studies therefore evaluate both direction and return as complementary views of the same forecasting task \cite{sezer2020dlreview, fischer2018lstm, nelson2017lstm}.

The corresponding input spaces also vary. Traditional models rely on historical prices, returns, volume, and hand-crafted technical indicators. These variables summarise past market behaviour, but they do not describe the semantic content of new information. A news article about profits, financing, or governance may precede a price move even when the numerical history alone is ambiguous. This observation motivates the use of exogenous text signals in addition to market data. The broader forecasting literature increasingly treats stock prediction as a multimodal problem rather than a purely numerical one \cite{atsalakis2009survey, sezer2020dlreview, fischer2018lstm}.

\section{Class Imbalance and Evaluation Sensitivity}

Directional forecasting introduces a statistical issue that is easy to underestimate: class imbalance. Upward and downward moves are rarely distributed evenly across a financial sample. Persistent market drift, low-volatility periods, and sparse firm-level event coverage can all push the labels toward one class. In that setting, raw accuracy overstates model quality because a classifier can appear successful while still failing on the minority class. This thesis therefore reports balanced accuracy, F1, and confusion-matrix statistics alongside headline accuracy, in line with the evaluation caution emphasised in recent forecasting studies \cite{sezer2020dlreview, liapis2023deep, shi2021deeplearning}.

The problem becomes sharper in sentiment-based firm-level studies. News coverage is uneven across firms, and longer forecast horizons can create clusters of same-direction labels. A model then learns the dominant class more easily than it learns a robust relation between sentiment and future movement. Table~\ref{tab:imbalance_background} summarises why imbalance matters and which safeguards are needed before a result can be interpreted as economically meaningful.

\begin{table}[ht]
\centering
\caption{Why class imbalance is a methodological concern in directional financial forecasting.}
\label{tab:imbalance_background}
\small
\renewcommand{\arraystretch}{1.15}
\begin{tabular}{>{\raggedright\arraybackslash}p{3.0cm}>{\raggedright\arraybackslash}p{4.1cm}>{\raggedright\arraybackslash}p{3.3cm}>{\raggedright\arraybackslash}p{3.1cm}}
\toprule
Imbalance source & Financial forecasting manifestation & Risk if ignored & Appropriate safeguard \\
\midrule
Market drift & One direction dominates over a sample period, especially in short binary up/down labels & Raw accuracy overstates true discrimination & Balanced accuracy and confusion-matrix reporting \\
Sparse firm coverage & Company panels with few relevant articles generate unstable minority-class observations & High F1 can be driven by majority-class prediction & Per-ticker evaluation and coverage analysis \\
Longer target horizons & Multi-day returns can cluster on one side during sustained moves & Thresholds become biased toward the dominant class & Validation-only threshold tuning and horizon-specific reporting \\
Event concentration & Market-moving news arrives in bursts rather than uniformly across dates & A model can memorise event periods instead of learning a stable relation & Chronological splits with embargo and walk-forward validation \\
\bottomrule
\end{tabular}
\end{table}

\section{Textual Information and Sentiment in Finance}

The finance literature treats text as an information-bearing object rather than as unstructured noise. Early work examined whether media tone, message-board language, or corporate disclosures explain subsequent market behaviour. Negative media language correlated with market activity and price pressure \cite{tetlock2007media}. Message-board sentiment also contained information about volume and volatility, although the signal quality depended on aggregation and filtering \cite{das2007yahoo}. Finance-specific dictionaries later showed that general-purpose sentiment lexicons can misclassify domain vocabulary, which is why terms such as ``liability'' or ``capital'' require financial context to be interpreted correctly \cite{loughran2011liability, kearney2014survey}.

These findings motivate a distinction between sentiment source and sentiment quality. Not every text stream has the same predictive value. Financial news is more structured than social-media text and is more likely to contain firm-linked events. Social media is timelier and larger in volume, but it is also noisier and harder to map to a specific asset. Survey work in financial sentiment analysis emphasises that source choice, content analysis method, and aggregation strategy jointly determine whether textual sentiment becomes a useful explanatory variable \cite{kearney2014survey, nguyen2015social, bollen2011twitter}.

\section{From Lexicons to Deep Financial NLP}

Text-based market prediction evolved from sparse lexical features toward distributed and contextual representations. Early financial news systems used bag-of-words, noun phrases, named entities, and support-vector models. The AZFinText system used breaking financial news and intra-day stock data to predict short-term price movement, showing that structured textual features support directional prediction \cite{schumaker2009azfin}. A later context-capturing study reported improvements over simpler textual baselines for stock-price prediction from financial news \cite{hagenau2013automated}. These systems demonstrated that the representation of text strongly influences the usefulness of the downstream forecasting model.

Deep learning extended this line of work by learning latent semantic features directly from text. Event-driven approaches represented news events in dense form and modelled how short-term and long-term event effects influence price movement \cite{ding2015event}. Other models fused news with recurrent architectures, attention mechanisms, or hybrid text-price representations \cite{hu2018chaotic, mohan2019news, xu2018stock}. These studies moved the field from manually designed textual cues toward learned semantic encodings. The same studies made clear that the meaning of a financial article cannot be reduced to isolated polarity words alone.

The broader NLP shift toward transformer-based models is also central to this thesis. The transformer architecture replaced recurrent sequence processing with self-attention, which enables stronger contextual representations at scale \cite{vaswani2017attention}. The Bidirectional Encoder Representations from Transformers (BERT) model then showed how pre-trained bidirectional language models can be adapted to downstream tasks with minimal architectural change \cite{devlin2019bert}. Domain-specialised variants followed, including FinBERT for English financial text \cite{araci2019finbert}. This progression matters because sentiment extraction quality increasingly depends on pre-training domain and language rather than on task-specific feature engineering alone.

\section{Arabic and Bilingual Financial Text Processing}

Arabic financial NLP introduces additional complexity. Arabic is morphologically rich, orthographically variable, and distributed across Modern Standard Arabic, dialectal Arabic, and specialised domain language. These properties make direct transfer from English sentiment models unreliable. Arabic-specific pre-training is therefore important for downstream classification quality. AraBERT demonstrated strong gains over multilingual baselines on Arabic NLP tasks \cite{antoun2020arabert}. ARBERT and MARBERT showed that specialised Arabic pre-training outperforms larger multilingual models on a broad benchmark of Arabic understanding tasks \cite{abdulmageed2021arbert}.

Bilingual settings introduce a further modelling decision: whether to force all text through a single multilingual model or to route each language through a more suitable encoder. Recent Arabic model studies indicate that language variant and task proximity are more important than raw pre-training size \cite{inoue2021camelbert, abdulmageed2021arbert, antoun2020arabert}. For a bilingual financial-news pipeline, these findings justify treating language routing as a methodological design choice rather than a preprocessing detail.

\section{Chronological Evaluation and Leakage}

The final background concept is evaluation discipline. Time-series forecasting is unusually vulnerable to leakage because future information can enter the feature space in subtle ways. Leakage may come from same-day article alignment, overlapping target windows, hyperparameter tuning on test periods, or refitting logic that uses observations whose targets are not yet known at prediction time. Reported performance can therefore be overstated even when the implementation looks reasonable on the surface \cite{liapis2023deep, dahal2023comparative}.

This concern is especially relevant for text-based forecasting. News has timestamps, trading occurs on exchange calendars, and forecast horizons span multiple future observations. A method that ignores those interactions can produce optimistic but invalid results. Chronological train-validation-test splits, embargo regions between adjacent blocks, and walk-forward evaluation are not optional refinements. They are part of the problem definition itself \cite{liapis2023deep, sezer2020dlreview}. This thesis adopts that view and treats leakage control as a core experimental requirement rather than as a final diagnostic step.

\section{Chapter Summary}

This chapter established the conceptual foundations for the thesis. Stock forecasting is difficult because market predictability is limited and unstable. Textual information attracts attention because it describes events that numerical histories do not represent directly. The usefulness of that information depends on representation quality, label balance, language suitability, and evaluation discipline. These background concepts motivate the more specific state-of-the-art review in the next chapter, which examines how prior work operationalised these ideas and where the present thesis contributes.
